\section{Discussion}

The experiments support a simple operational reading. A utility model for
auxiliary contexts is a ranking device, and its scores should not be treated
as measurements. Once that is accepted, the useful question is not how
accurate the scores are but how much of the routing decision can be certified.
Budget calibration answers that question with a number the operator chooses:
at two harmful selections per hundred decisions, the router acts on roughly
one decision in twelve; at five per hundred, roughly one in six. The remaining
decisions fall back to the anchor, which is a valid prediction rather than a
failure.

The comparison with interval certification is the sharpest result. Coverage
guarantees are the standard tool for distribution-free reliability, and they
are the wrong tool here, not because they are invalid but because they answer
a question about the error distribution rather than about the accepted set. In
a regime where the estimation error dominates the estimated quantity, the
lower bound is negative almost everywhere and the router abstains on almost
everything. The two rules share the same calibration data and the same
guarantee type; only the object being controlled differs, and the coverage
differs by a factor of four to nine.

Transfer is the practical caveat. A threshold calibrated on five benchmarks
transferred to the sixth with a shift margin below the nominal budget, and
across independent training runs the same held. The margin grows when the
budget is loosened, because a larger accepted set reaches further into the
score distribution where calibration states are sparse. Proposition
\ref{prop:shift} makes this quantitative, and the recommendation that follows
is to calibrate at the strictest budget the application can tolerate and to
re-estimate the threshold whenever the deployment distribution changes
materially.

Two limitations bound the claims. First, we control the mass of harmful
selections per decision, not the loss of the final prediction: a sequence of
acceptable selections can still produce a worse forecast than the anchor, and
the end-to-end study is the setting in which that residual risk is measured.
Second, the harmful event is defined against the true marginal utility, which
the router never observes; the analysis is therefore an evaluation-time
statement about a deployment-time rule, and its validity rests on the
exchangeability of calibration and deployment states. When that assumption
fails, the shift margin of Proposition~\ref{prop:shift} degrades with the
discrepancy, and the guarantee is only as good as the calibration sample.
